\section{Supply and establishment}
\label{sec:supplyestab}

Section~\ref{sec:doubleedge} ended by separating rescue into two requirements: a resistance mutation has to arise (supply) before extinction, and once it does, it has to grow to dominance in the drug-exposed edge, where antibiotic selection makes resistance an advantage (establishment). To understand where the dose-dependent flip in Figure~\ref{fig:doubleedge} comes from, we measure each requirement separately. The comparison between biofilm and well-mixed then splits into two pieces: how many mutations each population produces, and what fraction of those mutations establish in the edge.

Before splitting rescue into its two pieces, we first check that the bar-plot flip is a continuous feature of the dose response, not an artifact of sampling at two specific doses. Figure~\ref{fig:dosecurves} Panel A shows the three populations across the full dose sweep. The active-core biofilm here is drawn at our default $\bCore = 0.2$, which sets the core's birth and death rate to one-fifth of the edge birth rate $\bEdge = 1$; this is considerably less active than the $\bCore = 0.8$ used for Figure~\ref{fig:doubleedge}, where the core divides and dies at four-fifths of the edge rate. The active-core biofilm crosses the well-mixed reference smoothly near $\dEdge \approx 1.4$ and stays above it at every higher dose; the dormant-core biofilm ($\bCore = 0$, no metabolism in the interior) sits below the well-mixed reference at every dose. Panel B plots the same data as differences from the well-mixed reference: the active biofilm's $\Delta P(\mathrm{rescue})$ crosses zero near $\dEdge \approx 1.4$, and the dormant biofilm's $\Delta P(\mathrm{rescue})$ stays below zero throughout the sweep. Both bar-plot patterns therefore persist across the full dose range, and the crossover is robust to the level of core activity: even at the more modest $\bCore = 0.2$, an active interior is enough to flip the comparison at high dose.

\begin{figure}[h!]
\centering
\includegraphics[width=0.75\textwidth]{dose_curves_rescue.png}
\caption{Probability of rescue versus dose for the three populations: well-mixed reference (open circles), biofilm with a dormant core ($\bCore = 0$, open squares), and biofilm with an active core ($\bCore = 0.2$, filled triangles). Panel A: absolute $P(\mathrm{rescue})$, log $y$-axis. Panel B: difference relative to the well-mixed reference, $\Delta P(\mathrm{rescue}) = P_{\mathrm{BF}} - P_{\mathrm{WM}}$, with the dotted horizontal line at zero marking the no-difference baseline. Each point averages $5{,}000$ independent simulations. \emph{Parameters}: $\Nzero = 1000$, $l = 100$, $\mu = 10^{-4}$.}
\label{fig:dosecurves}
\end{figure}

Figure~\ref{fig:decomposition} Panel A reports the first piece of the decomposition: the difference in mutation count $\Delta S = S_{\mathrm{BF}} - S_{\mathrm{WM}}$ between the biofilm and the well-mixed reference, where $S$ is the mean number of R mutations a population produces per simulation run. Two observations. First, the dormant-core biofilm sits essentially on the $y = 0$ reference line at every dose: the dormant biofilm and the well-mixed reference produce essentially the same number of mutations, even though the dormant biofilm lives longer. We return to why that extra lifetime does not translate into extra mutations in a later section. Second, the active-core biofilm sits above the $y = 0$ line at every dose: an actively dividing interior produces more mutations than the well-mixed reference at every dose. The number of extra mutations the active biofilm produces over the well-mixed reference shrinks with dose, from roughly $1.5$ extra mutations per run at $\dEdge = 1.05$ down to $\sim 0.03$ at $\dEdge = 2.8$, because the absolute number of mutations any population produces shrinks at higher dose (shorter lifetimes give fewer cell divisions during which a mutation can occur). In ratio terms, however, the active biofilm produces about $1.7\times$ as many mutations as the well-mixed reference at every dose. Notably, the number of extra mutations the active biofilm produces does not track where the biofilm wins: the largest excess in mutation count is at low dose, precisely where the active biofilm rescues less often than the well-mixed reference (Figure~\ref{fig:dosecurves} Panel B); the excess shrinks toward zero at high dose, where the biofilm actually rescues more often. Where the active biofilm has the most additional mutations available, it converts them into the fewest rescue events.

\begin{figure}[h!]
\centering
\includegraphics[width=0.95\textwidth]{decomposition.png}
\caption{Decomposing the biofilm-to-well-mixed comparison into the two ingredients rescue requires. Panel A: difference in mutation supply $\Delta S = S_{\mathrm{BF}} - S_{\mathrm{WM}}$, where $S$ is the mean number of R mutations a population produces over a simulation run. Symlog $y$-axis, with the dotted line at $y = 0$ marking equal supply. Panel B: difference in per-mutation establishment fraction for edge-born mutations, $\Delta P(\mathrm{est}) = P_{\mathrm{est},\,\mathrm{BF}} - P_{\mathrm{est},\,\mathrm{WM}}$, where $P(\mathrm{est})$ is the fraction of edge-born R lineages whose populations reach the establishment threshold $k_{\mathrm{est}} = 3$ cells before going extinct. Linear $y$-axis, with the dotted line at $y = 0$ marking equal establishment. Both panels show two biofilm conditions across the dose sweep: dormant core ($\bCore = 0$, open squares) and active core ($\bCore = 0.2$, filled triangles). Each point averages $5{,}000$ independent runs. \emph{Parameters}: $\Nzero = 1000$, $l = 100$, $\mu = 10^{-4}$.}
\label{fig:decomposition}
\end{figure}

Figure~\ref{fig:decomposition} Panel B reports the second piece: the difference in per-mutation establishment between biofilm and well-mixed, $\Delta P(\mathrm{est}) = P_{\mathrm{est},\,\mathrm{BF}} - P_{\mathrm{est},\,\mathrm{WM}}$, where $P(\mathrm{est})$ is the fraction of edge-born R lineages whose populations reach the establishment threshold $k_{\mathrm{est}} = 3$ cells before going extinct. Both biofilm curves sit below the $y = 0$ reference line at every dose: a given edge-born R lineage is consistently less likely to grow past the early-extinction bottleneck inside a biofilm than in the well-mixed reference, by roughly $0.05$ to $0.10$ in absolute probability across most of the dose range. This reduced establishment applies to both dormant and active biofilms and persists across the entire sweep. We return to why edge-born R lineages are less likely to escape early extinction inside a biofilm in a later section.
